\DocumentMetadata{
  pdfversion=2.0,
  pdfstandard=ua-2,
  lang=en-US,
  tagging=on,
  tagging-setup={math/setup=mathml-SE,tabsorder=structure}
}
\documentclass[11pt,letterpaper]{article}
\usepackage[margin=0.68in,includefoot]{geometry}
\usepackage{newcomputermodern}
\usepackage{amsmath,amscd,mathtools}
\providecommand{\square}{\mdlgwhtsquare}
\usepackage[hidelinks]{hyperref}
\hypersetup{
  pdftitle={Logic PhD Exam, August 2019},
  pdfauthor={University of Florida Department of Mathematics},
  pdfsubject={Graduate mathematics examination},
  pdfdisplaydoctitle=true,
  bookmarksopen=true
}
\setcounter{secnumdepth}{0}
\setlength{\parindent}{0pt}
\setlength{\parskip}{0.28em}
\setlength{\emergencystretch}{3em}
\setlength{\leftmargini}{2.15em}
\setlength{\leftmarginii}{2.65em}
\setlength{\leftmarginiii}{3em}
\renewcommand{\labelenumi}{\textbf{\theenumi.}}
\pagestyle{empty}
\raggedbottom
\allowdisplaybreaks
\newenvironment{examproblems}[1]{%
  \begin{enumerate}%
  \setcounter{enumi}{#1}%
  \setlength{\topsep}{0.35em}%
  \setlength{\partopsep}{0pt}%
  \setlength{\itemsep}{0.48em}%
  \setlength{\parsep}{0.18em}%
}{\end{enumerate}}
\newenvironment{examsubparts}{%
  \begin{enumerate}%
  \setlength{\topsep}{0.18em}%
  \setlength{\partopsep}{0pt}%
  \setlength{\itemsep}{0.16em}%
  \setlength{\parsep}{0.1em}%
}{\end{enumerate}}
\newenvironment{examinstructions}{%
  \par\begingroup\itshape
}{%
  \par\endgroup\vspace{0.18em}
}
\makeatletter
\renewcommand\section{\@startsection{section}{1}{0pt}%
  {-1.25ex plus -.25ex minus -.1ex}%
  {0.55ex plus .1ex}%
  {\normalfont\large\bfseries}}
\renewcommand{\@maketitle}{%
  \newpage\null\vskip -1.2em%
  {\footnotesize\bfseries
    \href{https://www.ufl.edu/}{University of Florida}, \href{https://math.ufl.edu/}{Department of Mathematics}\par}%
  \vskip 0.18em%
  \tagstructbegin{tag=Title}%
    {\LARGE\bfseries\href{https://gma.math.ufl.edu/exams/logic-phd/}{Logic PhD Exam}, August 2019\par}%
  \tagstructend%
  \vskip 0.35em%
  \tagmcbegin{artifact}%
    \hrule height 1pt%
  \tagmcend%
  \vskip 0.55em%
}
\makeatother
\title{Logic PhD Exam, August 2019}
\author{}
\date{}
\begin{document}
\maketitle
\thispagestyle{empty}
\begin{examinstructions}
Solve 5 problems of the following; at least one from each section.
\end{examinstructions}
\section{A. Set Theory.}
\begin{examproblems}{0}
\item
Sketch the construction of any model of ZFC and the negation of the Continuum Hypothesis.
\item
Define the class of \(\Delta_0\) formulas. Prove carefully that if \(\phi(x)\) is a \(\Delta_0\) formula with one free variable and~\(M\) is any transitive set and \(x \in M\), then \(\phi(x)\) is equivalent to \(M \models \phi(x)\).
\item
Prove that if~\(X\) is a Polish space and \(A_0, A_1 \subset X\) are disjoint analytic sets, then there are disjoint Borel sets \(B_0, B_1 \subset X\) such that \(A_0 \subset B_0\) and \(A_1 \subset B_1\).
\end{examproblems}
\section{B. Computability.}
\begin{examproblems}{0}
\item
Define \(0'\) (zero jump) and show that it is a computably enumerable set which is not computable.
\item
Define the many-one reducibility and sketch the proof that there is a many-one degree strictly between~\(0\) and \(0'\).
\item
What is Gödel’s diagonalization lemma? Prove the lemma.
\end{examproblems}
\section{C. Model theory.}
\begin{examproblems}{0}
\item
State the downward Löwenheim–Skolem theorem and prove it.
\item
State the Łoś’ theorem and prove it.
\item
Find a complete theory with exactly one infinite countable model up to isomorphism. Prove this property of the theory.
\end{examproblems}
\end{document}
